% Figure polypeps_reading_order.
% Owner: none (note-only figure).
% Used by: gapped-2d-notes PP Sec. 2 (Lemma 2, step 1: the chain-rule order).
% Source: arXiv OpenAI2026PolynomialPEPS, sections/02-information.tex:180-223
% (proof of lem:cell-cmi: the colouring of the exterior, eq:info-tile-cost, the
% chain rule listing F_0 in creation order, the exceptional sites, then F_1 in
% reverse creation order, and the purity step I(U:C|Z) = I(U:C|(UCZ)^c)).
% Ink-to-index: dots are lattice sites; there are no tensors or open indices.
% Columns are blocks of sites in listing order, not geometry: B (cols 1-2, the
% forbidden-0 exterior, teal), the family-0 tiles U^0_1, U^0_2, U^0_3 in creation
% order (cols 3-5, teal), the exceptional sites (col 6, brown), the family-1 tiles
% in reverse creation order U^1_3, U^1_2, U^1_1 (cols 7-9, olive), the rest of
% the forbidden-1 exterior F_1 \ C (cols 10-11, olive) and the far set C (cols
% 12-13, blue). Rows 3-5 are sites; rows 1-2 and 6 are blank (bracket and label
% rows; inkless anchors only). The top bracket spans the conditioning order: B,
% then what the chain rule lists, in the order of the arrow. Two tiles are drawn
% red instead of their family hue (the tile U of the chain-rule term; one red
% contour, not a red halo around the family contour, per shared-spec 4.1 S3; the
% superscript keeps the family): for U = U^0_2 (col 4) the conditioning set Z is B and
% the sites listed before U, cols 1-3 (bottom bracket); for U = U^1_2 (col 8) the
% purity of Omega swaps Z = cols 1-7 for (UCZ)^c = cols 9-11 (bottom bracket),
% which holds only forbidden-1 sites and U^1_1, a family-1 tile created earlier
% than U^1_2. So each conditioning set is the tile's own forbidden colour plus
% earlier tiles of its family, which is eq:info-tile-cost.
% Not to scale: tiles are unions of O(u)-chambers of every generation, and there
% are C(1+log s)^K tiles and exceptional sites; three tiles per family are drawn.
% Boundary signature: empty.
% Label macros: none.
% Fit: one column (\tnshow[5.5mm]).
% Related: tn2d_two_family_purity (the purity flip for whole families, in the
% geometry of a half-plane; this figure is the order of the chain rule inside one
% cell, with the conditioning set of one term per family); polypeps_cell_holes
% (how the tiles are built); \chainglyph (the order glyph of the note).
\begin{center}
    \tenkzkernel
    \begin{tenkz}[rows={wire,wire,ket,ket,ket,wire}, cols=13, bonds=none]
        \tndeclare{species}{fam0}{hue=source:teal}
        \tndeclare{species}{fam1}{hue=source:olive}
        \tndeclare{species}{exceptional}{hue=source:brown}
        \tn[at=(1,1), skin=none]{} \tn[at=(1,9), skin=none]{}
        \tn[at=(6,1), skin=none]{} \tn[at=(6,3), skin=none]{}
        \tn[at=(6,9), skin=none]{} \tn[at=(6,11), skin=none]{}
        \tnmark[form=enclosure, species=fam0, tint, label pos=n]{(3,1) .. (5,2)}{$B$}
        \tnmark[form=enclosure, species=fam0, tint, label pos=n]{(3,3) .. (5,3)}{$U^0_1$}
        \tnmark[form=enclosure, species=selected, tint, label pos=n]{(3,4) .. (5,4)}{$U^0_2$}
        \tnmark[form=enclosure, species=fam0, tint, label pos=n]{(3,5) .. (5,5)}{$U^0_3$}
        \tnmark[form=enclosure, species=exceptional, tint, label pos=n]{(3,6) .. (5,6)}{$x$}
        \tnmark[form=enclosure, species=fam1, tint, label pos=n]{(3,7) .. (5,7)}{$U^1_3$}
        \tnmark[form=enclosure, species=selected, tint, label pos=n]{(3,8) .. (5,8)}{$U^1_2$}
        \tnmark[form=enclosure, species=fam1, tint, label pos=n]{(3,9) .. (5,9)}{$U^1_1$}
        \tnmark[form=enclosure, species=fam1, tint, label pos=n]{(3,10) .. (5,11)}{$F_1{\setminus}C$}
        \tnmark[form=enclosure, species=secondary, tint, label pos=n]{(3,12) .. (5,13)}{$C$}
        \tnmark[form=bracket, label pos=n]{(1,1) .. (1,9)}{$\text{listing order}\;\longrightarrow$}
        \tnmark[form=bracket, label pos=s]{(6,1) .. (6,3)}{$Z$ for $U^0_2$}
        \tnmark[form=bracket, label pos=s]{(6,9) .. (6,11)}{$(UCZ)^c$ for $U^1_2$}
    \end{tenkz}
\end{center}
